\section{The source Gaussian supremum and derivative estimates}
\label{sec:sr-source-gaussian-sup}
We use the source's Euclidean normalization
\(H(z,w)=\sum_i z_i\overline{w_i}\), Lebesgue measure on
\(\mathbb C^N\), and Gaussian \(e^{-q\pi\|z\|^2}\).
The square of the \(L^2\)-norm is denoted \(L_2\), as in the source.

\subsection{The actual lattice and its covolume}
Let \(Y=\mathrm{GT}(\iota)\), \(N=|\iota|\).
Its period lattice \(L\subset\mathbb C^N\) is generated by a
real basis \(b_1,\ldots,b_{2N}\).
The alternating form \(\operatorname{Im}H\) is integral on \(L\),
and every \(\mathbb Z\)-linear functional on \(L\) is pairing
with a point of \(L\). The supplied factor \(k(a)\) has
\(|k(a)|=1\) for \(a\in L\).
An actual level-\(q\) section is an entire \(F\) satisfying
\[
 F(z+a)=k(a)\exp\{q\pi H(z,a)+q\pi H(a,a)/2\}F(z).
 \tag{GS1}\label{eq:sr-source-GT-section}
\]
Its norm density and the source normalized integral are
\[
 h_F(z)=|F(z)|^2e^{-q\pi\|z\|^2},\qquad
 L_2(F)={1\over\operatorname{vol}D}\int_Dh_F,\qquad
 \lambda_1=\inf_{0\ne a\in L}\|a\|,
 \tag{GS2}\label{eq:sr-source-GT-norm}
\]
where \(D\) is the half-open basis parallelepiped.
Equation (GS1) gives \(h_F(z+a)=h_F(z)\) directly.

The source first proves \(\operatorname{vol}D\le1\).
Let \(B\) be the real basis matrix in orthonormal coordinates
and \(J\) the orthogonal complex structure representing
\(\operatorname{Im}H\), with \(|\det J|=1\).
The alternating Gram matrix is \(G=B^tJ^tB\).
For each lattice coordinate functional, perfection supplies a
lattice vector representing it.
Writing these vectors in the given integral basis produces
an integer matrix \(C\) with \(CG=I\).
Thus \(\det C\ne0\), \(|\det C|\ge1\), and
\[
                  1=|\det C|\,|\det B|^2,\qquad
                           \operatorname{vol}D=|\det B|\le1 .
 \tag{GS3}\label{eq:sr-source-GT-volume}
\]
This follows the source
\srcid{LemmaSA.lemSG_inv}, \srcid{lemS_detM_le},
and \srcid{lemS_vol_le}.
In particular the integral over \(D\) is at most \(L_2(F)\).
For \(N>0\), discreteness of the basis lattice and the existence
of a nonzero basis vector give \(0<\lambda_1<\infty\).

\subsection{The local Gaussian mean}
For an entire \(G:\mathbb C^N\to\mathbb C\), \(N>0\),
\[
 |G(0)|^2\le {2\over\pi^N}
       \int_{\|z\|\le\sqrt{2N}}|G(z)|^2e^{-\|z\|^2}\,dz .
 \tag{GS4}\label{eq:sr-source-gaussian-ball-mean}
\]
Here is the source's mean argument.
For fixed \(z\), the one-variable Cauchy mean and
Cauchy--Schwarz give
\(|G(0)|^2\le\int_{|u|=1}|G(uz)|^2\,du\), with probability
measure on the circle.
Integrating on a radial ball and using the measure-preserving
rotations \(z\mapsto uz\) proves
\[
 |G(0)|^2\int_{\|z\|\le\sqrt{2N}}e^{-\|z\|^2}
       \le\int_{\|z\|\le\sqrt{2N}}|G(z)|^2e^{-\|z\|^2}.
\]
The total Gaussian mass is \(\pi^N\), and its second moment is
\(N\pi^N\). On \(\|z\|>\sqrt{2N}\),
\(1\le\|z\|^2/(2N)\), so that region has mass at most
\(\pi^N/2\).
This proves (GS4).
The source verifies these two moments through the radial
integrals
\(\int_0^\infty r^ke^{-r^2}dr=\Gamma((k+1)/2)/2\).
For \(N=0\) the ball has its single point with mass \(1\);
(GS4) remains true.
These are the steps in \srcid{RSubMean.sq_norm_zero_le_gauss_ball}.

For \(q\ge1\), \(v_0\in\mathbb C^N\), put
\(s=\sqrt{q\pi}\), \(w=z/s\), and
\[
 G(z)=F(v_0+w)\exp\{-q\pi H(w,v_0)\}.
\]
Then
\[
 |G(z)|^2e^{-\|z\|^2}
              =e^{q\pi\|v_0\|^2}h_F(v_0+w).
\]
Changing real \(2N\)-dimensional volume by \(s^{-2N}\) in
(GS4) gives the exact local estimate
\[
 h_F(v_0)\le2q^N
       \int_{\|u-v_0\|\le\sqrt{2N}/\sqrt{q\pi}}h_F(u)\,du .
 \tag{GS5}\label{eq:sr-source-local-Gaussian}
\]
This is \srcid{LemmaSB.lemS_local}, with its actual square-root
and \(\pi\) normalization.

\subsection{Packing and the global estimate}
A finite \(\delta\)-separated subset of a ball of radius \(\rho\)
in \(\mathbb C^N\) has cardinality at most
\[
                         (1+4\rho/\delta)^{2N}.
 \tag{GS6}\label{eq:sr-source-complex-packing}
\]
To match the source's constants, when \(\rho<\delta/2\)
there is at most one point.
Otherwise place disjoint open balls of radius \(\delta/2\)
around the points and compare their real \(2N\)-dimensional
volumes with the containing ball of radius
\(\rho+\delta/2\). This gives
\((1+2\rho/\delta)^{2N}\); the source's
\(\mathrm{PackingCount}\) bound \( (3\max(\delta,\rho)/\delta)^{2N}\)
also implies (GS6) in the stated range.
This is the conversion in \srcid{LemmaS.lemS_pack}.

Let \(\rho=\sqrt{2N}/\sqrt{q\pi}\) and
\(K=(1+4\rho/\lambda_1)^{2N}\).
For any fixed \(u\in D\), the lattice translates \(u+a\) lying
in the closed ball of radius \(\rho\) around \(v_0\) form a
\(\lambda_1\)-separated set.
Apply (GS6) to each finite subset; the nonnegative sum of their
indicators is therefore at most \(K\).
Unfolding Lebesgue measure over \(D+L\), Tonelli and the
periodicity of \(h_F\) yield
\[
 \int_{\|u-v_0\|\le\rho}h_F(u)\,du
   \le K\int_Dh_F(u)\,du\le K L_2(F).
 \tag{GS7}\label{eq:sr-source-ball-cover}
\]
The density is continuous and the closure of \(D\) is compact,
so the latter integral is finite.
This justifies the conversion from the nonnegative integrals,
including the finite-measure requirements.
For \(N=0\) the same estimate is the one-point calculation.

\begin{srproposition}\label{prop:sr-source-gaussian-sup}
For every actual \(F\in Y.\mathrm{Sec}(q)\), \(q\ge1\),
\[
 |F(v)|^2e^{-q\pi\|v\|^2}
 \le 2q^N
       \left(1+{4\sqrt{2N}\over\sqrt{q\pi}\lambda_1}\right)^{2N}
                     L_2(F)\qquad(v\in\mathbb C^N).
 \tag{GS8}\label{eq:sr-source-gaussian-sup}
\]
For \(N=0\), the parenthesized factor is interpreted by the
separate zero-dimensional calculation.
\end{srproposition}
\begin{proof}
Combine (GS5) and (GS7).
This is exactly \srcid{LemmaS.lemS_main}, then
\srcid{gSup} and \srcid{GlobalTheory.sup_le_l2_principal_torus}.
\end{proof}

\subsection{The source repeated Cauchy estimate}
Let \(F\) be entire on \(\mathbb C^N\), and suppose
\[
                   |F(z)|^2e^{-\pi\|z\|^2}\le S
                                      \quad\text{for every }z.
\]
Necessarily \(S\ge0\).
For a list of \(d\) coordinate directions,
\[
 |D_{i_1}\cdots D_{i_d}F(0)|
                   \le\sqrt S\,(\pi e d)^{d/2}.
 \tag{GS9}\label{eq:sr-source-Gaussian-Cauchy}
\]
For \(d=0\) the last factor is \(1\).
The source proves this by repeated one-variable Cauchy estimates.
If \(|F|\le A\) whenever \(\sum_i|z_i-x_i|\le d\rho\), induction
on the list gives the derivative bound \(A/\rho^d\).
In the induction step apply the one-variable Cauchy derivative
bound on the circle of radius \(\rho\) in the next coordinate;
each remaining ball stays in the original \(d\rho\)-ball.

For \(d>0\), take \(\rho=(\pi d)^{-1/2}\).
On that \(d\rho\)-ball about zero,
\(\|z\|^2\le(\sum_i|z_i|)^2\le d/\pi\).
The Gaussian hypothesis gives \(|F(z)|\le\sqrt S\,e^{d/2}\).
Consequently the repeated bound is
\(\sqrt S\,e^{d/2}(\pi d)^{d/2}\), which is (GS9).
This is the actual \srcid{RCauchy.iter_step_bound},
\srcid{cauchy_gauss_le} and \srcid{GCauchy.gCauchy} argument.
For a homogeneous polynomial \(P\) of degree \(d\), finite
linearity gives
\[
 |P(\partial)F(0)|
     \le \left(\sum_\alpha|P_\alpha|\right)
                           \sqrt S\,(\pi e d)^{d/2}.
 \tag{GS10}\label{eq:sr-source-polynomial-Cauchy}
\]
The source polynomial derivative is the sum over monomials of
the iterated derivative along their actual repeated coordinate
list; no factorial or divided-power convention is inserted.

\paragraph{The tensor-period minimum used with this estimate.}
Let \(T\) be the source \(O_E\)-linear positive torus, and
\(\Lambda\subset E^n\) the actual positive self-dual Hermitian
coefficient lattice.
The source short-vector supplier gives an \(E\)-spanning family
\(v_j\in\Lambda\) satisfying
\[
 \sum_k|\tau(v_{j,k})|^2\le
               \{2 f^{3/2}n^2\mathrm{rd}_E\}^2
                         \quad\text{for all }j,\tau .
 \tag{GS11}\label{eq:sr-source-tensor-short-vectors}
\]
For \(\gamma\) in the actual tensor lattice, contraction by
\(v_j\) lies in the original period lattice.
If \(\gamma\ne0\), some such contraction is nonzero:
otherwise linearity and the spanning property make all
contractions zero, hence every coordinate of \(\gamma\) zero.
Cauchy--Schwarz gives
\[
 \|\operatorname{contr}_{v_j}\gamma\|
     \le2 f^{3/2}n^2\mathrm{rd}_E\,\|\operatorname{flat}\gamma\|.
\]
Thus, if every nonzero original period has length at least
\(\lambda>0\), every nonzero tensor period has length at least
\[
                  \lambda'={\lambda\over
                                     2f^{3/2}n^2\mathrm{rd}_E}.
 \tag{GS12}\label{eq:sr-source-tensor-period-minimum}
\]
For \(n=0\) the tensor lattice is zero and the source takes
\(\lambda'=1\).
These are the exact contraction steps in
\srcid{LeMulNormFlatOfSelfDual.le_mul_norm_flat_of_selfDual}
and \srcid{NormalisedThetaValueBound.exists_min_period}.
The short-vector construction (GS11) is its separately retained
arithmetic lattice supplier.

\paragraph{Frozen-source ledger.}
All sources use \srcid{alpoge-source-7cbf7a0c}.
In
\srcid{src/ReciprocityTheory/ReductiveGroups/ArithmeticGeometry/SquareIntegrability/},
the relevant files and blobs are
\begin{itemize}
\item \srcid{GCauchyStatement.lean},
\srcid{e53576b249c74c71fcb3af8c66b5d3111df34e04};
\item \srcid{LemmaSA.lean},
\srcid{dbdc66e846263d475c90d6c0c9ef0e1c580b845a};
\item \srcid{RSubMean.lean},
\srcid{d6dda1e32c137194deb3cf921537bf7c21e6efe4};
\item \srcid{LemmaSB.lean},
\srcid{948ae7e2d3bc2acbcf2b3a97ad9b4ceae6fe1609};
\item \srcid{LemmaS.lean},
\srcid{1950e4726b3b5ce278491239f53fcab26d875218};
\item \srcid{RCauchy.lean},
\srcid{07fa9a67f8ac8ff589a18c3691ec38b76f963408};
\item \srcid{GCauchy.lean},
\srcid{22dc8c28fccd0c9fad68fd47686cd6e0a01cef93}.
\end{itemize}
The realization of the packing supplier by disjoint ball volumes
is \srcid{src/Development/AbelianVarietyTheory/Proofs/CuspFormEvaluations/SupLeL2PrincipalTorus.lean},
blob \srcid{89e09ea66f30995d745920d4a3f50db3921618ba}.
The tensor contraction proof in the same directory is
\srcid{LeMulNormFlatOfSelfDual.lean},
blob \srcid{0800466bdb169dc1ddd430812943ca9c4f8f093d};
the actual normalized tensor-period bound is
\srcid{NormalisedThetaValueBound.lean},
blob \srcid{689a12045c213d4d11e73f6660d5ee8a8739f684}.
The Gaussian statements above are proved for a supplied actual
\(\mathrm{GT}\); they do not construct the model's analytic
fibre or identify an algebraic theta reading with its sections.
